\section{Definitions of symbols}

This appendix distinguishes the mathematical relations named in the
package. The descriptions are conventions for these notes, not claims
that every author uses the same symbols.

\begin{description}[leftmargin=2.8em,style=nextline]
\item[$\in$] Membership: $x\in A$ says that $x$ is an element of $A$.
\item[$\subseteq$] Inclusion: $A\subseteq B$ allows $A=B$.
\item[$\subsetneq$] Proper inclusion: $A\subsetneq B$ requires
  $A\subseteq B$ and $A\ne B$. The bare symbol $\subset$ is used
  differently by different authors; prefer $\emphpropcontainedin$ when
  strictness matters.
\item[$\vdash$] Syntactic derivability: $\Gamma\proves\varphi$
  states that a proof of $\varphi$ from $\Gamma$ exists in the chosen
  deductive system.
\item[$\vDash$] Semantic consequence: $\Gamma\entails\varphi$
  states that every relevant structure satisfying $\Gamma$ also
  satisfies $\varphi$.
\item[$\models$] Satisfaction: $\Str\satisfies\varphi$ says that a
  particular structure $\Str$ satisfies $\varphi$. The glyph closely
  resembles or coincides with the previous one; the custom command
  records which relation is intended.
\item[$\Vdash$] Forcing: $w\forces\varphi$ uses the forcing relation
  appropriate to the stated Kripke or poset semantics.
\item[$\succeq$] Accessibility in the chosen Kripke orientation,
  written \verb|\kripkesucceq|. The legacy \verb|\Acc| is available
  with the \verb|[kripke]| option. Peano successor remains $S$ and
  is not assigned this relation symbol.
\item[$\geqslant$] Slanted greater-than-or-equal, written
  \verb|\gecurly|. This is an order glyph, not the Kripke relation.
\item[$\langle A,R\rangle$] A displayed structure or tuple.
  \verb|\struc{A,R}| and \verb|\tuple{A,R}| print the same brackets
  while recording the intended role in source.
\end{description}

\subsection{A two-item check}
Decide whether each statement is true or false:
\[
  \text{(1) }\emptyset\in\emptyset,
  \qquad
  \text{(2) }\emptyset\subset\emptyset.
\]
The first is false: the empty set has no elements. The second depends
on the author's convention for $\subset$: it is true if equality is
allowed and false if strict inclusion is intended. Written without
ambiguity, $\emptyset\subseteq\emptyset$ is true and
$\emptyset\subsetneq\emptyset$ is false.
